% ============================================================================
% Fichier  : partie3/P03_C08_famille_iso27000.tex
% Titre    : La Famille ISO/IEC 27000
% Auteur   : MINKA MI NGUIDJOÏ Thierry Emmanuel
% Version  : 1.0 -- novembre 2026
% Corrigés : corriges/Corr_C08.tex
% ============================================================================
% Vérifié sur iso.org le 06/11/2026 : ISO/IEC 27002:2022, édition 3, publiée
% en février 2022 (remplace l'édition 2013 et ses rectificatifs) ; ISO/IEC
% 27001:2022 et son amendement 1:2024 ; ISO/IEC 27000:2018 en vigueur.
% Textes sous licence : articles et mesures PARAPHRASÉS, jamais cités ;
% intitulés des mesures traduits par le cours. Le dossier Normes ne contient
% que l'édition 2013 : écart signalé dans le texte.
% Structure de l'édition 2022 (93 mesures ; 37/8/14/34 ; 11 mesures
% nouvelles ; 5 attributs), date de fin de transition (31 octobre 2025),
% objet de l'amendement 2024 (changement climatique) : d'après des sources
% publiques concordantes, à confirmer sur le texte par l'auteur.
% Faits du cas (bible) : certificat de juin 2025 (CERTIQUAL AFRIQUE, 3 ans,
% centre de données seul) ; périmètre étendu approuvé par le conseil
% (Ex. 1.4) ; résultats de l'audit de mars 2028 (T1 à T8).
% Faits nouveaux (reportés dans la bible) : renouvellement du certificat en
% juin 2028 sur le périmètre étendu ; extrait de la déclaration
% d'applicabilité (10 lignes) ; exclusion de 8.28 contredite par les
% interfaces développées par la DSI ; revue croisée 3/3/3.
% ============================================================================

\chapter{La Famille ISO/IEC 27000}
\label{chap:iso27000}
\chaptermark{La famille ISO/IEC 27000}

\epigraph{%
    « Une norme ne dit pas ce que vous faites.\\
    Elle dit ce que vous devrez prouver. »%
}{Principe du cours}

\niveauChapitre{Fondamental}{%
    Maîtrise du Chapitre~\ref{chap:iso27005}%
}

\begin{boxobjectifs}
\begin{itemize}
    \item Situer les principales normes de la famille ISO/IEC~27000 et leur
          rôle.
    \item Présenter les exigences de l'ISO/IEC~27001:2022, articles~4 à~10,
          et le cycle de certification.
    \item Utiliser l'ISO/IEC~27002:2022~: 93 mesures, quatre thèmes,
          attributs, et l'écart avec l'édition de 2013.
    \item Rédiger une déclaration d'applicabilité et justifier une
          exclusion.
    \item Confronter une déclaration d'applicabilité aux preuves d'audit.
    \item Situer les ISO/IEC~27004, 27006 et 27007.
\end{itemize}
\end{boxobjectifs}

\bigskip

% ============================================================================
\section*{Avant de commencer : le renouvellement}
% ============================================================================

Douala, avril 2028. Le certificat ISO/IEC~27001 de la BANQUE
SANAGA\footnote{Comme dans tout le cours, la BANQUE SANAGA, l'organisme
CERTIQUAL AFRIQUE, les dates et les montants sont fictifs~; les normes
citées sont exactes et paraphrasées, leur texte étant protégé.}, obtenu en
juin 2025 pour trois ans, arrive à échéance. Le conseil a approuvé
l'extension du système de management aux 38 agences~; l'audit de
renouvellement de CERTIQUAL AFRIQUE portera donc, pour la première fois,
sur toute la banque.

Mme~ABENA doit remettre à l'organisme certificateur une pièce maîtresse~:
la \emph{déclaration d'applicabilité}, qui dit, mesure par mesure, ce que
la banque applique et pourquoi. La version de 2025 tenait en quelques pages,
sur le seul centre de données. Celle de 2028 doit couvrir les agences, au
moment même où l'audit interne vient d'y relever trois non-conformités et
deux inconnues (Chapitre~\ref{chap:conclure}).

\medskip

Ce chapitre ouvre la Partie~III, consacrée au cadre normatif et
juridique. Il présente la famille de normes la plus utilisée en sécurité
de l'information, puis la manière de rédiger, et de vérifier, une
déclaration d'applicabilité.

\begin{boxalert}
\textbf{Violation typique.} Rédiger la déclaration d'applicabilité en
cochant « mise en œuvre » pour toutes les mesures, à partir des
déclarations des services.

\smallskip
\textbf{Règle~:} la déclaration d'applicabilité est un document
\emph{opposable}~: l'organisme certificateur et les auditeurs la
confrontent aux preuves. Chaque « mise en œuvre » doit pouvoir être
démontrée~; chaque exclusion doit être justifiée par l'appréciation des
risques.
\end{boxalert}

\separateur

% ============================================================================
\section{Une Famille de Normes}
\label{sec:c08-famille}
% ============================================================================

La famille ISO/IEC~27000 est publiée par le sous-comité~27 du comité
technique commun de l'ISO et de la CEI (JTC~1/SC~27). Elle compte des
dizaines de documents~; une seule norme est certifiable, l'ISO/IEC~27001.

\begin{tableaucours}{L{2.6cm}L{5.4cm}L{4.2cm}}{%
    \tableauHeader{Norme} &
    \tableauHeader{Objet} &
    \tableauHeader{Dans ce cours}
}
ISO/IEC~27000 & Vue d'ensemble et vocabulaire~\cite{iso27000_2018} &
Définitions (propriétaire de risque) \\
ISO/IEC~27001 & Exigences du SMSI~; seule norme certifiable de la
famille~\cite{iso27001_2022} & Ce chapitre, Chapitre~\ref{chap:smsi} \\
ISO/IEC~27002 & Recommandations pour les mesures de
sécurité~\cite{iso27002_2022} & Ce chapitre \\
ISO/IEC~27004 & Surveillance, mesure, analyse et évaluation~\cite{iso27004_2016}
& Chapitre~\ref{chap:smsi} \\
ISO/IEC~27005 & Gestion des risques~\cite{iso27005_2022} &
Chapitre~\ref{chap:iso27005} \\
ISO/IEC~27006 & Exigences pour les organismes qui auditent et certifient
les SMSI & Ce chapitre \\
ISO/IEC~27007 & Lignes directrices pour l'audit des SMSI~\cite{iso27007_2020}
& Chapitres~\ref{chap:audit} et~\ref{chap:smsi} \\
ISO/IEC~27014 & Gouvernance de la sécurité~\cite{iso27014_2020} &
Chapitre~\ref{chap:grc} \\
\end{tableaucours}
\vspace{-0.8em}
\begin{center}{\small\itshape Tableau~8.1. Les principales normes de la
famille ISO/IEC~27000}\end{center}

D'autres normes de la famille déclinent les mesures pour un secteur ou un
sujet~: les services en nuage, la gestion des incidents, la protection de
la vie privée. Le Chapitre~\ref{chap:emergents} y revient pour le nuage.

% ============================================================================
\section{L'ISO/IEC 27001:2022}
\label{sec:c08-27001}
% ============================================================================

\subsection{Une structure commune}

L'ISO/IEC~27001 suit la structure commune à toutes les normes de systèmes
de management de l'ISO (qualité, environnement, continuité)~: dix
articles, dont les sept derniers contiennent les exigences. Une
organisation qui a déjà un système de management de la qualité retrouve
donc le même plan.

\begin{tableaucours}{cL{3.0cm}L{8.4cm}}{%
    \tableauHeader{Art.} &
    \tableauHeader{Titre} &
    \tableauHeader{Ce que la norme exige, en substance}
}
4 & Contexte de l'organisation & Comprendre les enjeux internes et
externes, les parties intéressées et leurs exigences~; fixer le
périmètre~; établir le SMSI et ses processus \\
5 & Leadership & Engagement de la direction~; politique de sécurité~;
rôles, responsabilités et autorités \\
6 & Planification & Actions face aux risques et opportunités~:
appréciation et traitement des risques, déclaration d'applicabilité~;
objectifs de sécurité~; planification des modifications \\
7 & Support & Ressources~; compétences~; sensibilisation~;
communication~; informations documentées \\
8 & Fonctionnement & Planifier et maîtriser les opérations~; réaliser les
appréciations et les traitements des risques prévus \\
9 & Évaluation des performances & Surveillance et mesure~; audit
interne~; revue de direction \\
10 & Amélioration & Amélioration continue~; non-conformités et actions
correctives \\
\end{tableaucours}
\vspace{-0.8em}
\begin{center}{\small\itshape Tableau~8.2. Les articles~4 à~10 de
l'ISO/IEC~27001:2022 (paraphrasés)}\end{center}

Les Chapitres~\ref{chap:grc} à~\ref{chap:conclure} ont déjà parcouru
l'essentiel de ces exigences sans les nommer~: le périmètre et les parties
intéressées (article~4), le rôle de la direction (article~5),
l'appréciation et le traitement des risques (articles~6 et~8), l'audit
interne (article~9). La norme ne prescrit pas \emph{comment} faire~; elle
dit ce que l'organisation doit faire et pouvoir prouver.

\subsection{Ce qui a changé en 2022}

L'édition de 2022 remplace celle de 2013. Les principaux changements~:
\begin{itemize}
    \item l'annexe~A, liste des mesures de référence, est alignée sur la
          nouvelle ISO/IEC~27002:2022 (93 mesures au lieu de 114)~;
    \item une exigence nouvelle impose de planifier les modifications du
          SMSI (article~6.3)~;
    \item le SMSI doit identifier ses processus et leurs interactions
          (article~4.4).
\end{itemize}
Un amendement publié en 2024 demande en outre à l'organisation de
déterminer si le changement climatique est un enjeu pertinent pour son
système de management. Les certificats délivrés sur l'édition de 2013
devaient avoir migré au plus tard le 31~octobre 2025. Le certificat de la
BANQUE SANAGA, délivré en juin 2025 sur l'édition de 2022, n'était pas
concerné.

\begin{boxinfo}
\textbf{Le dossier du cours et l'édition 2013.} Les normes archivées dans
les ressources du cours sont celles de 2013. Le cours enseigne l'édition
de 2022~; les numéros de mesures cités ici sont ceux de 2022. Un étudiant
qui consulte un document de 2013 trouvera d'autres numéros (de A.5 à A.18)
et quatorze domaines au lieu de quatre thèmes.
\end{boxinfo}

\subsection{Le cycle de certification}

\begin{tableaucours}{L{3.4cm}L{8.8cm}}{%
    \tableauHeader{Étape} &
    \tableauHeader{Contenu}
}
Audit initial, étape~1 & Revue documentaire et de l'état de préparation~:
périmètre, politique, appréciation des risques, déclaration
d'applicabilité \\
Audit initial, étape~2 & Évaluation de la mise en œuvre et de
l'efficacité du SMSI sur le terrain~; décision de certification \\
Audits de surveillance & Au moins un par an, pendant les trois ans du
certificat, sur une partie du SMSI \\
Audit de renouvellement & Avant l'échéance des trois ans, sur l'ensemble
du SMSI \\
\end{tableaucours}
\vspace{-0.8em}
\begin{center}{\small\itshape Tableau~8.3. Le cycle de certification}\end{center}

L'organisme certificateur doit lui-même satisfaire aux exigences de
l'ISO/IEC~27006, et être accrédité par un organisme d'accréditation. La
certification est volontaire~: elle ne remplace ni l'audit interne, ni
l'audit de sécurité réglementaire (Chapitre~\ref{chap:audit}).

% ============================================================================
\section{L'ISO/IEC 27002:2022}
\label{sec:c08-27002}
% ============================================================================

\subsection{93 mesures, quatre thèmes}

L'ISO/IEC~27002 n'est pas certifiable~: elle donne, pour chaque mesure, un
objectif et des recommandations de mise en œuvre~\cite{iso27002_2022}.
L'édition de 2022 réorganise les mesures en quatre thèmes.

\begin{tableaucours}{L{3.4cm}cL{7.0cm}}{%
    \tableauHeader{Thème} &
    \tableauHeader{Mesures} &
    \tableauHeader{Exemples (intitulés traduits par le cours)}
}
Organisationnelles (5.1 à 5.37) & 37 & Politiques, rôles, renseignement
sur les menaces (5.7), relations avec les fournisseurs (5.19, 5.20),
préparation des TIC pour la continuité (5.30) \\
Liées aux personnes (6.1 à 6.8) & 8 & Sélection, sensibilisation et
formation (6.3), signalement des événements \\
Physiques (7.1 à 7.14) & 14 & Périmètres, surveillance de la sécurité
physique (7.4), matériels \\
Technologiques (8.1 à 8.34) & 34 & Droits d'accès privilégiés (8.2),
authentification sécurisée (8.5), vulnérabilités techniques (8.8),
sauvegarde (8.13), journalisation (8.15), cloisonnement des réseaux
(8.22), codage sécurisé (8.28) \\
\end{tableaucours}
\vspace{-0.8em}
\begin{center}{\small\itshape Tableau~8.4. Les quatre thèmes de
l'ISO/IEC~27002:2022}\end{center}

Onze mesures sont nouvelles, dont le renseignement sur les menaces (5.7),
la sécurité des services en nuage (5.23), la préparation des TIC pour la
continuité (5.30), la surveillance de la sécurité physique (7.4), la gestion
de la configuration (8.9), la prévention des fuites de données (8.12), les
activités de surveillance (8.16), le filtrage web (8.23) et le codage
sécurisé (8.28). Les autres mesures de 2013 ont été regroupées ou
reformulées.

\subsection{Les attributs}

L'édition de 2022 associe à chaque mesure des \termecle{attributs}\index{attributs
de mesure}, qui permettent de trier les mesures selon plusieurs vues~:
\begin{itemize}
    \item le \emph{type de mesure}~: préventive, détective, corrective~;
    \item les \emph{propriétés de sécurité}~: confidentialité, intégrité,
          disponibilité~;
    \item les \emph{concepts de cybersécurité}~: identifier, protéger,
          détecter, répondre, rétablir (les fonctions historiques du cadre
          du NIST, Chapitre~\ref{chap:grc})~;
    \item les \emph{capacités opérationnelles} (gestion des identités,
          continuité, gestion des fournisseurs, etc.)~;
    \item les \emph{domaines de sécurité}~: gouvernance et écosystème,
          protection, défense, résilience.
\end{itemize}

\begin{boxinfo}
\textbf{À quoi servent les attributs~?} À vérifier l'équilibre d'un
dispositif. Les mesures de la BANQUE SANAGA en 2025 étaient presque toutes
\emph{préventives}~; les mesures \emph{détectives} (journalisation,
surveillance) et \emph{correctives} (sauvegardes, mode dégradé) manquaient
en agence. Trier la déclaration d'applicabilité par type de mesure l'aurait
montré. C'est aussi ce que faisaient apparaître les étapes du scénario
(Chapitre~\ref{chap:risque}).
\end{boxinfo}

% ============================================================================
\section{La Déclaration d'Applicabilité}
\label{sec:c08-dda}
% ============================================================================

\subsection{Ce qu'elle contient}

\begin{boxdef}
\termecle{Déclaration d'applicabilité}\index{déclaration d'applicabilité}
(DdA)~: document exigé par l'ISO/IEC~27001 (article~6.1.3) qui recense les
mesures nécessaires au traitement des risques, justifie leur inclusion,
indique si elles sont mises en œuvre, et justifie l'exclusion de toute
mesure de l'annexe~A~\cite{iso27001_2022}.
\end{boxdef}

La déclaration fait le pont entre l'appréciation des risques
(Chapitre~\ref{chap:iso27005}) et la liste des mesures. Une mesure y figure
parce qu'un risque l'exige, ou parce qu'une obligation légale ou
contractuelle l'impose. Une mesure en est exclue parce qu'aucun risque ni
aucune obligation ne la rend nécessaire, et non parce qu'elle est difficile.

\subsection{Un extrait}

\begin{boxscenario}{Opération SANAGA. La déclaration d'applicabilité de 2028 (extrait)}
\begin{tableaucours}{cL{3.6cm}cL{4.4cm}c}{%
    \tableauHeader{Mesure} &
    \tableauHeader{Intitulé (traduit)} &
    \tableauHeader{Incl.} &
    \tableauHeader{Justification} &
    \tableauHeader{Mise en œuvre}
}
5.20 & Sécurité dans les accords avec les fournisseurs & Oui & R-DC-02,
R-AG-02 & Oui \\
5.30 & Préparation des TIC pour la continuité & Oui & R-AG-01, récupération
& Oui \\
6.3 & Sensibilisation et formation & Oui & R-AG-01, entrée & Oui \\
8.2 & Droits d'accès privilégiés & Oui & R-AG-01, propagation & Oui \\
8.5 & Authentification sécurisée & Oui & R-AG-01, entrée & Oui \\
8.8 & Gestion des vulnérabilités techniques & Oui & R-AG-01, R-AG-02 & Oui \\
8.13 & Sauvegarde des informations & Oui & R-AG-01, récupération & Oui \\
8.15 & Journalisation & Oui & R-DC-03, loi n°~2010/012, art.~29 & Oui \\
8.22 & Cloisonnement des réseaux & Oui & R-AG-01, propagation & Oui \\
8.28 & Codage sécurisé & Non & « Pas de développement interne~: le
système central est un progiciel » & -- \\
\end{tableaucours}

\smallskip
{\small La première version remise par les directions, en avril 2028,
porte « Oui » dans la colonne « mise en œuvre » pour les neuf mesures
incluses.}
\end{boxscenario}

\subsection{Confronter la déclaration aux preuves}

La RSSI a appris la leçon de février~: avant d'envoyer la déclaration,
elle la confronte aux résultats de l'audit de mars~2028
(Chapitres~\ref{chap:audit} et~\ref{chap:conclure}).

\begin{boxlecture}{Les neuf « mise en œuvre » face aux preuves}
\textbf{Constatation.}
\begin{itemize}
    \item \OUI{}~: 8.5 (second facteur, T1), 8.13 (sauvegardes, T3), 8.22
          (cloisonnement, T2)~;
    \item \NON{}~: 5.30 (mode dégradé, T4), 8.2 (administration locale,
          T7), 8.15 (journalisation incomplète pendant la période, T6)~;
    \item \NEK{}~: 6.3 (sensibilisation, T5), 5.20 (le contrat des GAB n'a
          pas de droit d'audit, T8), 8.8 (vulnérabilités~: GAB non vérifiés,
          postes non testés).
\end{itemize}
3~\OUI{}, 3~\NON{}, 3~\NEK{}~:
\[
\mu = U = \tfrac{1}{3}~;\quad
H = \log_2 3 \approx 1{,}585~\text{bit}.
\]

\textbf{Analyse} (confiance élevée). La déclaration affirmait neuf mesures
sur neuf~; les preuves en confirment trois. La dispersion est maximale~:
tout ce que l'on peut dire de cette déclaration, c'est qu'elle ne
renseigne pas sur l'état réel. La version envoyée à CERTIQUAL AFRIQUE est
corrigée~: « partielle » pour 5.30, 8.2 et 8.15, avec le plan d'actions~;
« à vérifier » pour 6.3, 5.20 et 8.8, avec la date de la vérification
prévue.

\medskip
\hfill\lecture{0{,}33}{0{,}33}{1{,}58}
\end{boxlecture}

\subsection{L'exclusion de 8.28}

L'exclusion du codage sécurisé paraît solide~: la banque n'écrit pas son
système central. Mais la DSI développe, depuis 2026, les interfaces entre
le système central, l'application de paiement et le portail en ligne.
L'exclusion est donc contredite par un fait~: sa base est \NON{}. La mesure
est réintégrée, limitée aux développements d'interfaces.

\begin{boxprocedure}
\textbf{Vérifier une exclusion.} Le papier de travail qui examine une
exclusion de la déclaration d'applicabilité vérifie~:
\begin{itemize}
    \item que la justification renvoie à l'appréciation des risques (aucun
          risque ne rend la mesure nécessaire), et non à une difficulté de
          mise en œuvre~;
    \item qu'aucune obligation légale ou contractuelle n'impose la
          mesure~;
    \item que les faits invoqués sont exacts (ici~: existe-t-il vraiment
          aucun développement interne~?)~;
    \item que l'exclusion est approuvée par la direction, et datée.
\end{itemize}
\end{boxprocedure}

% ============================================================================
\section{Mesurer et Auditer le SMSI}
\label{sec:c08-autres}
% ============================================================================

Deux normes de la famille accompagnent le fonctionnement du SMSI.
L'ISO/IEC~27004 aide à construire des indicateurs qui répondent à
l'exigence de surveillance et de mesure de l'article~9.1~\cite{iso27004_2016}.
L'ISO/IEC~27007 adapte à l'audit d'un SMSI les lignes directrices générales
de l'ISO~19011~\cite{iso27007_2020}. Le Chapitre~\ref{chap:smsi} met en
œuvre l'une et l'autre~: les indicateurs de la RSSI, puis l'audit interne du
SMSI élargi et la préparation de la revue de direction.

\begin{boxinfo}
\textbf{La norme et la loi.} L'article~26 de la loi n°~2010/012 impose aux
exploitants des systèmes d'information de se doter de « systèmes
normalisés » de gestion des risques (Chapitre~\ref{chap:iso27005})~\cite{loi2010012}.
Un SMSI conforme à l'ISO/IEC~27001 en est une forme reconnue
internationalement~; la loi n'en fait pas une obligation, et le certificat
ne dispense pas de l'audit de sécurité réglementaire des articles~13 et~32.
Le Chapitre~\ref{chap:droit} précise l'articulation.
\end{boxinfo}

% ============================================================================
\section{Ce que la Norme demande à chaque Métier}
\label{sec:c08-metiers}
% ============================================================================

\begin{boxmetier}{Consultant en mise en œuvre de SMSI}
\textbf{Une déclaration défendable}
\begin{enumerate}
    \item Partir de l'appréciation des risques, pas de la liste de
          l'annexe~A.
    \item Pour chaque mesure incluse, citer le risque ou l'obligation qui
          l'exige.
    \item Pour chaque « mise en œuvre », désigner la preuve qui le
          démontre.
    \item Justifier chaque exclusion par le risque, et en vérifier les
          faits.
\end{enumerate}
\end{boxmetier}

\begin{boxmetier}{Auditeur en cybersécurité}
\textbf{Lire la déclaration comme un ensemble d'assertions}

Chaque ligne « mise en œuvre » est une assertion de l'audité~; l'auditeur
de certification la teste, comme Mme~EKAMBI a testé le plan de traitement.
Il prélève en priorité les mesures nouvelles de 2022, les exclusions, et
les mesures qui portent les risques les plus élevés.
\end{boxmetier}

\begin{boxmetier}{RSSI}
\textbf{Ne rien déclarer que l'on ne peut prouver}

La RSSI préfère une déclaration honnête, avec des « partielle » et des
« à vérifier », à une déclaration parfaite qui ne résisterait pas à
l'étape~2 de l'audit. Un organisme certificateur pardonne une mesure en
cours~; il ne pardonne pas une déclaration fausse.
\end{boxmetier}

\begin{boxmetier}{Responsable GRC}
\textbf{Relier la déclaration au registre}

La colonne « justification » de la déclaration renvoie aux références du
registre des risques (R-AG-01, R-DC-02\ldots). Le responsable GRC vérifie
que tout risque intolérable a ses mesures dans la déclaration, et que toute
mesure de la déclaration a un risque ou une obligation.
\end{boxmetier}

\separateur

\begin{boxresume}
\textbf{Ce qu'il faut retenir de ce chapitre~:}
\begin{itemize}
    \item \textbf{La famille ISO/IEC~27000} compte une seule norme
          certifiable, l'ISO/IEC~27001~; les autres l'outillent.
    \item \textbf{L'ISO/IEC~27001:2022} suit la structure commune des
          systèmes de management~: articles~4 à~10, du contexte à
          l'amélioration.
    \item \textbf{L'ISO/IEC~27002:2022} compte 93 mesures en quatre thèmes
          (37 organisationnelles, 8 liées aux personnes, 14 physiques, 34
          technologiques), avec des attributs.
    \item \textbf{La déclaration d'applicabilité} relie les risques aux
          mesures~; chaque « mise en œuvre » se prouve, chaque exclusion se
          justifie par le risque.
    \item \textbf{Confrontée aux preuves}, la déclaration de la BANQUE
          SANAGA passe de 9 sur 9 à 3 \OUI{}, 3 \NON{}, 3 \NEK{}.
    \item \textbf{La certification} est volontaire et triennale~; elle ne
          remplace ni l'audit interne, ni l'audit réglementaire.
\end{itemize}
\end{boxresume}

\begin{boxinfo}
\textbf{Points de contrôle.} Vous avez compris ce chapitre si vous savez~:
\begin{itemize}
    \item dire dans quel article de l'ISO/IEC~27001 se trouve une exigence~;
    \item placer une mesure dans son thème et lui associer des attributs~;
    \item rédiger une ligne de déclaration d'applicabilité avec sa
          justification~;
    \item contester une exclusion par les faits.
\end{itemize}
\end{boxinfo}

% ============================================================================
\section*{Exercices}
% ============================================================================

\begin{boxexo}[title={\ding{45}\quad Exercice~8.1. Quel article~? \niveaubase}]
Indiquez l'article de l'ISO/IEC~27001:2022 (4 à 10) dont relève chacune des
activités suivantes.
\begin{enumerate}[(a)]
    \item Le conseil approuve l'extension du périmètre aux agences.
    \item La RSSI rédige la déclaration d'applicabilité.
    \item L'audit interne conduit la mission de mars 2028.
    \item La DSI planifie le remplacement du collecteur de journaux sans
          interrompre la journalisation.
    \item La direction du réseau d'agences corrige le mode dégradé de Kribi
          après l'observation S4.2.
    \item La RSSI organise la formation de deux personnes par agence.
\end{enumerate}
\end{boxexo}

\begin{boxexo}[title={\ding{45}\quad Exercice~8.2. Thèmes et attributs \niveaumoyen}]
Pour chacune des mesures suivantes, donnez son thème et son type
(préventive, détective, corrective), puis dites à quelle étape du scénario
R-AG-01 elle se rattache.
\begin{enumerate}[(a)]
    \item 8.5, authentification sécurisée.
    \item 8.15, journalisation.
    \item 8.13, sauvegarde des informations.
    \item 6.3, sensibilisation et formation.
    \item 8.22, cloisonnement des réseaux.
    \item 5.30, préparation des TIC pour la continuité.
\end{enumerate}
Que montre la répartition par type de mesure~?
\end{boxexo}

\begin{boxexo}[title={\ding{45}\quad Exercice~8.3. Quatre exclusions \niveaumoyen}]
Une banque de la sous-région exclut de sa déclaration d'applicabilité les
mesures suivantes. Pour chacune, dites si l'exclusion est recevable, et
justifiez.
\begin{enumerate}[(a)]
    \item 5.23 (services en nuage)~: « Nous n'utilisons aucun service en
          nuage. » La messagerie est hébergée chez un fournisseur en nuage.
    \item 7.4 (surveillance de la sécurité physique)~: « Trop coûteux pour
          les petites agences. »
    \item 8.28 (codage sécurisé)~: « Aucun développement~; tous les logiciels
          sont des progiciels, paramétrés par l'éditeur. » L'inventaire le
          confirme.
    \item 8.12 (prévention des fuites de données)~: « Mesure couverte par la
          politique de classification. »
\end{enumerate}
\end{boxexo}

\begin{boxexo}[title={\ding{45}\quad Exercice~8.4. La déclaration de LiZbiZ \niveauexpert}]
LiZbiZ veut s'engager vers la certification ISO/IEC~27001 pour rassurer
PHARMACAM (Chapitres~\ref{chap:ebios} à~\ref{chap:audit}).
\begin{enumerate}[(a)]
    \item Proposez un périmètre défendable, en tenant compte des
          dépendances (éditeur, transporteurs, centres de données).
    \item Rédigez six lignes de déclaration d'applicabilité pour les mesures
          qui traitent le scénario SO1, avec leur justification.
    \item L'audit contractuel de PHARMACAM a donné~: comptes par défaut
          supprimés (\OUI{}), services chiffrés (\OUI{}), filtrage entre
          centres (\NON{}), second facteur de l'éditeur (\NEK{}, accès
          refusé), restauration observée (\OUI{}), plan de réponse
          (\NON{}). Remplissez la colonne «~mise en œuvre~», calculez $\mu$,
          $U$ et $H$, et dites si LiZbiZ peut demander l'audit initial.
    \item Quelle mesure de l'ISO/IEC~27002 traite la cause du \NEK{}, et que
          doit faire LiZbiZ avant l'étape~2~?
\end{enumerate}
\end{boxexo}

\bigskip

\begin{boxinfo}
\textbf{Pour aller plus loin.}
\begin{itemize}
    \item L'ISO/IEC~27001:2022~\cite{iso27001_2022}, articles~4 à~10 et
          annexe~A.
    \item L'ISO/IEC~27002:2022~\cite{iso27002_2022}, pour les recommandations
          de mise en œuvre de chaque mesure.
    \item L'ISO/IEC~27000~\cite{iso27000_2018}, pour le vocabulaire.
\end{itemize}
\end{boxinfo}

\begin{boxinfo}
\textbf{Transition.} L'ISO/IEC~27001 n'est pas le seul référentiel. La
banque doit aussi répondre au superviseur, et ses partenaires parlent
parfois NIST ou CIS. Comment passer d'un référentiel à l'autre sans tout
refaire~? C'est l'objet du Chapitre~\ref{chap:referentiels}.
\end{boxinfo}
